% ============================================================
% Chapter 2 — Section 2.6: Research Gap
% Draft for master thesis body (not the proposal)
% ============================================================
%
% Usage: paste into thesis main .tex, inside \section{Background and Related Work}
% Dependencies: \usepackage{booktabs}, \usepackage{multirow}, \usepackage{array}
%               (all already present in thesis_proposal.tex template)
% ============================================================

\subsection{Research Gap}
\label{sec:gap}

The preceding sections have reviewed GUI quality assessment, LLM-generated
interfaces, LLM-as-a-judge methodology, and multimodal web agents.
Table~\ref{tab:related-work} places the most relevant prior works along
seven dimensions that together characterise a comprehensive GUI quality
evaluation framework.
Taken together, three observations motivate the present thesis.

\paragraph{No prior work applies modern multimodal LLMs as multi-dimensional GUI judges.}
UIClip~\cite{uiclip} demonstrated that vision-language models can perform
pairwise UI comparison, but it relies on models that predate the current
generation of multimodal LLMs and limits evaluation to a binary preference.
The LLM-as-a-judge literature~\cite{llm_judge,li-etal-2025-generation}
has established systematic judge evaluation for text, yet this methodology
has not been transferred to the UI domain.

\paragraph{Requirement fidelity and rubric-based assessment are largely absent from GUI evaluation.}
Design2Code~\cite{design2code} evaluates how faithfully generated code
reproduces a visual reference, but does not assess richer quality dimensions
such as visual hierarchy, information clarity, or perceived usability.
The Designer Feedback study~\cite{designer_feedback} found that simple
pairwise rankings among evaluators agreed at only around 50\%, suggesting
that coarse binary judgements are insufficient for reliable GUI assessment.
A structured rubric grounded in established usability
theory~\cite{nielsen1990heuristic,iso9241-11} can provide finer-grained and
more reproducible quality signals.

\paragraph{Static visual judgements have never been validated against dynamic interaction outcomes.}
VisualWebArena~\cite{visualwebarena} and WebArena~\cite{webarena}
evaluate multimodal agents on goal-directed web tasks, but do not assess
the GUI quality of the interfaces themselves.
Whether a high static quality score predicts successful task completion
by a computer-use agent remains an open empirical question.
Answering this question requires a benchmark that keeps generated interfaces
executable so that static and dynamic evaluations can be performed on the
same set of items.

The present thesis addresses all three gaps simultaneously by combining
a two-track benchmark, a rubric-based human annotation protocol, a systematic
comparison of modern multimodal LLM judges with bias and reliability analysis,
and a dynamic evaluation module that links static scores to interaction outcomes.

% ----------------------------------------------------------------
\begin{table*}[t]
  \centering
  \caption{Comparison of related work across seven evaluation dimensions.
           \checkmark~= fully addressed;
           $\triangle$~= partially addressed;
           $\times$~= not addressed;
           ---~= not applicable to this line of work.}
  \label{tab:related-work}
  \small
  \setlength{\tabcolsep}{5pt}
  \begin{tabular}{lccccccc}
    \toprule
    \textbf{Work}
      & \begin{tabular}[c]{@{}c@{}}\textbf{Multi-}\\\textbf{dimensional}\\\textbf{rubric}\end{tabular}
      & \begin{tabular}[c]{@{}c@{}}\textbf{Modern}\\\textbf{multimodal}\\\textbf{LLM judge}\end{tabular}
      & \begin{tabular}[c]{@{}c@{}}\textbf{Bias \&}\\\textbf{reliability}\\\textbf{analysis}\end{tabular}
      & \begin{tabular}[c]{@{}c@{}}\textbf{Human}\\\textbf{reference}\\\textbf{labels}\end{tabular}
      & \begin{tabular}[c]{@{}c@{}}\textbf{Requirement}\\\textbf{fidelity}\end{tabular}
      & \begin{tabular}[c]{@{}c@{}}\textbf{Dynamic}\\\textbf{interaction}\\\textbf{eval}\end{tabular}
      & \begin{tabular}[c]{@{}c@{}}\textbf{Benchmark}\\\textbf{artifact}\end{tabular} \\
    \midrule
    UIClip~\cite{uiclip}
      & $\times$ & $\times$ & $\times$ & \checkmark & $\times$ & $\times$ & $\times$ \\
    Designer Feedback~\cite{designer_feedback}
      & $\times$ & $\times$ & $\times$ & \checkmark & $\triangle$ & $\times$ & $\times$ \\
    Design2Code~\cite{design2code}
      & $\triangle$ & \checkmark & $\times$ & \checkmark & \checkmark & $\times$ & \checkmark \\
    MT-Bench / Chatbot Arena~\cite{llm_judge}
      & $\triangle$ & \checkmark & $\triangle$ & \checkmark & --- & --- & $\triangle$ \\
    LLM-as-a-Judge survey~\cite{li-etal-2025-generation}
      & \checkmark & \checkmark & \checkmark & --- & --- & --- & $\times$ \\
    Position Bias Study~\cite{shi-etal-2025-judging}
      & $\times$ & \checkmark & \checkmark & \checkmark & --- & --- & $\times$ \\
    WebArena~\cite{webarena}
      & $\times$ & $\triangle$ & $\times$ & \checkmark & $\triangle$ & \checkmark & \checkmark \\
    VisualWebArena~\cite{visualwebarena}
      & $\times$ & \checkmark & $\times$ & \checkmark & $\triangle$ & \checkmark & \checkmark \\
    \midrule
    \textbf{This thesis}
      & \checkmark & \checkmark & \checkmark & \checkmark & \checkmark & \checkmark & \checkmark \\
    \bottomrule
  \end{tabular}
\end{table*}
% ----------------------------------------------------------------
